\documentclass[11pt,a4paper]{article}

% -- Core packages --
\usepackage{amsmath,amssymb}
\usepackage{geometry}
\usepackage{graphicx}
\usepackage{booktabs}
\usepackage{hyperref}
\usepackage{xcolor}
\usepackage{listings}
\usepackage{forest}         % for directory trees
\usepackage{tikz}
\usetikzlibrary{arrows.meta,positioning,shapes.geometric}

% -- Code listing style --
\lstset{
  basicstyle=\ttfamily\small,
  keywordstyle=\bfseries,
  commentstyle=\color{gray},
  stringstyle=\color{brown},
  breaklines=true,
  frame=none,
  language=C++,
  showstringspaces=false,
}

% -- coderef package (source-code cross-reference) --
\usepackage{coderef}
% Auto-generated by scan_coderefs.py — do not edit by hand.
% Regenerate with: python scan_coderefs.py <source_dir>

\coderefsetcount{axial_layer_state}{1}
\coderefsetentry{axial_layer_state}{1}{platform/FuelRodState.hpp}{15}

\coderefsetcount{cladding_creep_hook}{1}
\coderefsetentry{cladding_creep_hook}{1}{platform/CladdingMaterial.hpp}{41}

\coderefsetcount{cladding_iface}{1}
\coderefsetentry{cladding_iface}{1}{platform/CladdingMaterial.hpp}{12}

\coderefsetcount{fred_ida_solver_base}{1}
\coderefsetentry{fred_ida_solver_base}{1}{platform/FredIdaSolverBase.hpp}{6}

\coderefsetcount{fred_m_na_gap_behavior}{1}
\coderefsetentry{fred_m_na_gap_behavior}{1}{apps/fred\_m\_na/FredMNaGapBehavior.hpp}{19}

\coderefsetcount{fred_m_na_stress_strain}{1}
\coderefsetentry{fred_m_na_stress_strain}{1}{apps/fred\_m\_na/FredMNaStressStrain.hpp}{32}

\coderefsetcount{fred_m_na_stress_strain_residual}{1}
\coderefsetentry{fred_m_na_stress_strain_residual}{1}{apps/fred\_m\_na/FredMNaStressStrain.hpp}{43}

\coderefsetcount{fred_ox_global_residuals}{1}
\coderefsetentry{fred_ox_global_residuals}{1}{apps/fred\_ox/FredOxResiduals.hpp}{93}

\coderefsetcount{fred_ox_irradiation_hook}{1}
\coderefsetentry{fred_ox_irradiation_hook}{1}{apps/fred\_ox/FredOxResiduals.hpp}{96}

\coderefsetcount{fred_ox_layer_state}{1}
\coderefsetentry{fred_ox_layer_state}{1}{apps/fred\_ox/FredOxState.hpp}{10}

\coderefsetcount{fred_ox_residuals}{1}
\coderefsetentry{fred_ox_residuals}{1}{apps/fred\_ox/FredOxResiduals.hpp}{11}

\coderefsetcount{fred_ox_rod_state}{1}
\coderefsetentry{fred_ox_rod_state}{1}{apps/fred\_ox/FredOxState.hpp}{35}

\coderefsetcount{fred_ox_solver}{1}
\coderefsetentry{fred_ox_solver}{1}{apps/fred\_ox/FredOxSolver.hpp}{21}

\coderefsetcount{fred_rod_residuals}{1}
\coderefsetentry{fred_rod_residuals}{1}{apps/fred\_rod/FredRodResiduals.hpp}{11}

\coderefsetcount{fred_rod_run}{1}
\coderefsetentry{fred_rod_run}{1}{apps/fred\_rod/FredRodSolver.hpp}{35}

\coderefsetcount{fred_rod_solver}{1}
\coderefsetentry{fred_rod_solver}{1}{apps/fred\_rod/FredRodSolver.hpp}{15}

\coderefsetcount{fred_rod_thermal_hook}{1}
\coderefsetentry{fred_rod_thermal_hook}{1}{apps/fred\_rod/FredRodResiduals.hpp}{65}

\coderefsetcount{fred_rod_yvec_layout}{1}
\coderefsetentry{fred_rod_yvec_layout}{1}{apps/fred\_rod/FredRodResiduals.hpp}{19}

\coderefsetcount{fred_solver_base}{1}
\coderefsetentry{fred_solver_base}{1}{platform/FredSolverBase.hpp}{34}

\coderefsetcount{fuel_pellet_iface}{1}
\coderefsetentry{fuel_pellet_iface}{1}{platform/FuelPelletMaterial.hpp}{13}

\coderefsetcount{fuel_rod_geometry}{1}
\coderefsetentry{fuel_rod_geometry}{1}{platform/FuelRodGeometry.hpp}{27}

\coderefsetcount{gap_bond_selector}{1}
\coderefsetentry{gap_bond_selector}{1}{platform/GapMaterial.hpp}{32}

\coderefsetcount{gap_material_iface}{1}
\coderefsetentry{gap_material_iface}{1}{platform/GapMaterial.hpp}{21}

\coderefsetcount{gap_pressure_model}{1}
\coderefsetentry{gap_pressure_model}{1}{platform/GapPressureModel.hpp}{83}

\coderefsetcount{gap_state_manager}{1}
\coderefsetentry{gap_state_manager}{1}{platform/GapStateManager.hpp}{14}

\coderefsetcount{geometry_build}{1}
\coderefsetentry{geometry_build}{1}{platform/FuelRodGeometry.hpp}{73}

\coderefsetcount{heat_conduction_neq}{1}
\coderefsetentry{heat_conduction_neq}{1}{platform/HeatConduction.hpp}{38}

\coderefsetcount{heat_conduction_residual}{1}
\coderefsetentry{heat_conduction_residual}{1}{platform/HeatConduction.hpp}{51}

\coderefsetcount{rod_residuals_class}{1}
\coderefsetentry{rod_residuals_class}{1}{platform/RodResiduals.hpp}{18}

\coderefsetcount{rod_residuals_hooks}{1}
\coderefsetentry{rod_residuals_hooks}{1}{platform/RodResiduals.hpp}{133}

\coderefsetcount{rod_residuals_neq_layout}{1}
\coderefsetentry{rod_residuals_neq_layout}{1}{platform/RodResiduals.hpp}{192}

\coderefsetcount{rod_residuals_nvi}{1}
\coderefsetentry{rod_residuals_nvi}{1}{platform/RodResiduals.hpp}{99}

\coderefsetcount{rod_state}{1}
\coderefsetentry{rod_state}{1}{platform/FuelRodState.hpp}{134}

\coderefsetcount{stress_strain_neq}{1}
\coderefsetentry{stress_strain_neq}{1}{platform/StressStrain.hpp}{57}

\coderefsetcount{stress_strain_residual}{1}
\coderefsetentry{stress_strain_residual}{1}{platform/StressStrain.hpp}{64}

\coderefsetcount{update_geometry}{1}
\coderefsetentry{update_geometry}{1}{platform/FuelRodState.hpp}{85}
   % generated by scan_coderefs.py

% -- Layout --
\geometry{margin=2.5cm}

\hypersetup{
  colorlinks=true,
  linkcolor=blue!60!black,
  urlcolor=blue!70!black
}

% -- Title --
\title{%
  \textbf{FRED 2.0 Platform}\\[8pt]
  \large Developer Guide\\[4pt]
  \normalsize Code hierarchy and extension points for FRED-ROD and FRED-OX
}
\author{FRED Development Team\\
        \small Paul Scherrer Institut Advanced Nuclear Systems group}
\date{\today}

% ===========================================================================
\begin{document}
% ===========================================================================

\maketitle
\tableofcontents
\newpage

% ===========================================================================
\section{Introduction}
\label{sec:intro}
% ===========================================================================

This guide is intended for developers who want to understand the internal structure of the FRED~2.0 platform, add new material correlations, or build a new application on top of the platform layer.  It assumes familiarity with the physics: thermal conduction, thermo-elastic stress-strain, gap mechanics, and the SUNDIALS IDA time-integration scheme described in the Theory Manual.

\subsection*{Prerequisites}
\begin{itemize}
  \item C++17 toolchain (GCC~11 or later).
  \item SUNDIALS~6 (IDA, NVecSerial, SunMatrixDense, SunLinsolDense, Newton nonlinear solver) built as static archives.
  \item pybind11 and a Python 3.10+ conda environment for the Python bindings.
\end{itemize}

\subsection*{How to build}
From the repository root:
\begin{lstlisting}[language=bash]
make pymodule              # build fred.<ext>.so + fred_input.py + fred_output.py
make exe                   # build build/fred_rod.x (standalone C++ driver)
make install               # copy .so into the active conda env site-packages
\end{lstlisting}
The \texttt{CONDA\_ENV\_NAME} variable (default: \texttt{fred-dev}) and \texttt{SUNDIALS\_PREFIX} can be overridden on the command line.

% ===========================================================================
\section{Repository Layout}
\label{sec:layout}
% ===========================================================================

\begin{forest}
  for tree={font=\ttfamily\small, grow'=0, child anchor=west, parent
  anchor=south, anchor=west, calign=first,
  edge path={\noexpand\path [draw, \forestoption{edge}]
             (!u.south west) +(7.5pt,0) |- (.child anchor)\forestoption{edge label};},
  before typesetting nodes={if n=1{insert before={[,phantom]}}{}},
  fit=band, before computing xy={l=15pt}}
  [fred\_platform/
    [src/
      [platform/       , label=right:{abstract base classes and shared physics blocks}]
      [apps/
        [fred\_rod/    , label=right:{FRED-ROD: thermo-mechanics only}]
        [fred\_ox/     , label=right:{FRED-OX: + MOX irradiation physics}]
        [fred\_m\_na/  , label=right:{FRED-M-Na: + metallic fuel physics}]
      ]
      [bindings.cpp    , label=right:{pybind11 module root (\texttt{PYBIND11\_MODULE})}]
    ]
    [python/
      [fred\_input.py  , label=right:{Python validation wrappers}]
      [fred\_output.py , label=right:{HDF5 reader helpers}]
    ]
    [examples/         , label=right:{runnable Python examples with plots}]
    [doc/              , label=right:{theory manual and this guide}]
    [Makefile]
  ]
\end{forest}

\medskip
The \texttt{platform/} directory must not import anything from \texttt{apps/}. All physics shared across applications lives in \texttt{platform/}.  An application brings its own residual assembler, solver driver, extended state struct, and material classes.

% ===========================================================================
\section{Platform Layer}
\label{sec:platform}
% ===========================================================================

\subsection{Geometry and State Structures}
\label{sec:geom_state}

\paragraph{FuelRodGeometry.}
The struct \coderef{fuel_rod_geometry} holds the as-fabricated geometry of a single fuel rod: the number of radial nodes in the fuel (\texttt{nf}) and cladding (\texttt{nc}), the number of axial layers (\texttt{nz}), and the key radii \texttt{rfi0}, \texttt{rfo0}, \texttt{rci0}, \texttt{rco0}.  All lengths are in metres.

After setting the primary fields, the caller must invoke \coderef{geometry_build} to populate the derived arrays: uniform node spacings \texttt{drf0}/\texttt{drc0}, the full node-radius array \texttt{rad0[nf+nc]}, face-midpoint radii \texttt{rad\_half}, and control-volume areas \texttt{area0}.  These derived quantities are consumed by \texttt{HeatConduction} and \texttt{StressStrain}.

\paragraph{AxialLayerState.}
The struct \coderef{axial_layer_state} carries the complete runtime state for a single axial slice.  Fields are grouped into thermal, gap, fuel mechanics, and cladding mechanics. The deformed radii and layer heights are recomputed from the strain state via \coderef{update_geometry}, which returns \texttt{false} if any dimension goes negative (signalling the IDA solver to reject the step).

\paragraph{RodState.}
\coderef{rod_state} aggregates one \texttt{AxialLayerState} per axial layer together with the global gas-phase quantities (\texttt{gpres}, \texttt{mu0}, \texttt{fgrel}).  This is the state type used by FRED-ROD.  FRED-OX extends it with its own \texttt{FredOxRodState} (Section~\ref{sec:ox_state}).

\subsection{Material Interfaces}
\label{sec:materials}

All material properties are accessed through abstract C++ interfaces.  The platform layer depends only on these interfaces; concrete implementations live in the application directories.

\paragraph{FuelPelletMaterial.}
Declared at \coderef{fuel_pellet_iface}.  Required overrides:
\begin{itemize}
  \item \texttt{thermalConductivity(T)} — as-fabricated value in W/(m$\cdot$K).
  \item \texttt{heatCapacity(T)}, \texttt{thermalExpansionStrain(T)}.
  \item \texttt{youngsModulus(T, density)}, \texttt{poissonRatio()}.
  \item \texttt{referenceDensity()}, \texttt{theoreticalDensity()}, \texttt{meltingTemperature()}.
\end{itemize}
Irradiation-induced modifications (burnup-dependent conductivity, swelling correlations) are not part of this interface.  They are introduced by the application layer through the irradiation residual hook (Section~\ref{sec:rod_residuals}).

\paragraph{CladdingMaterial.}
Declared at \coderef{cladding_iface}.  Required overrides mirror those of \texttt{FuelPelletMaterial} plus \texttt{meyerHardness(T)} (used in the Ross--Stoute solid-contact conductance model).  The optional hook \coderef{cladding_creep_hook} returns the equivalent hoop creep strain rate $\dot{\varepsilon}_{c}$ [1/s].  It defaults to zero, making it a no-op for FRED-ROD cladding classes.  FRED-OX cladding classes (e.g.\ \texttt{FredOxT91}) override it with an irradiation-creep correlation that the irradiation residual block then integrates as an ODE.

\paragraph{GapMaterial.}
Declared at \coderef{gap_material_iface}.  The interface supports two bond types, selected by overriding \coderef{gap_bond_selector}:
\begin{description}
  \item[Gas bond (\texttt{isGasBond() = true}, default).] The platform calls \texttt{gapConductivity(T)} to obtain the gas mixture thermal conductivity $k$ [W/(m$\cdot$K)], then evaluates gap conductance from first principles using the Lanning--Hann temperature-jump model applied to the actual gap width, plus linearised radiation between concentric cylinders.
  \item[Liquid bond (\texttt{isGasBond() = false}).] The platform calls \texttt{gapConductivity(T)} to obtain $k$ [W/(m$\cdot$K)] and computes
    \begin{equation*}
      h = \frac{k(T)}{\max\!\bigl(\delta,\;\sqrt{r_\mathrm{uf}^2+r_\mathrm{uc}^2}\bigr)},
    \end{equation*}
    where $\delta$ is the current gap width and $r_\mathrm{uf}$, $r_\mathrm{uc}$ are the fuel and cladding surface roughnesses.  No gas-jump correction and no radiation term are applied.  After this calculation \texttt{clampConductance(h, gap\_closed)} is called; the default implementation returns $h$ unchanged, but subclasses may override it to enforce minimum/maximum bounds (e.g.\ sodium bond forces $h \geq 10^6$~W/(m$^2\cdot$K) at gap closure).
\end{description}

\subsection{Physics Blocks}
\label{sec:physics_blocks}

\paragraph{HeatConduction.}
Implements the finite-volume energy balance in cylindrical geometry for a single axial layer.  The equation count per layer is \coderef{heat_conduction_neq}: $4 + n_f + n_c$ equations (four global AEs for \texttt{qqv}, \texttt{gap}, \texttt{pfc}, \texttt{hgap}; $n_f$ fuel temperature ODEs; $n_c - 1$ inner cladding temperature ODEs; one outer cladding AE).

The method \coderef{heat_conduction_residual} fills the residual sub-vector for one layer.  It is called by \texttt{RodResiduals::computeResiduals} inside the per-layer loop and must not be called directly by application code.

\textbf{Coolant boundary condition modes.} The outer cladding node is always algebraic; its residual equation can take one of two forms (see the Theory Manual, Section~2.1 for the governing equations):
\begin{description}
  \item[Dirichlet (default, \texttt{isCladTempOuterBoundaryPrescribed = true}).] The outer cladding surface temperature is pinned directly to the supplied coolant temperature \texttt{T\_coolant}.  This is the mode used by FRED-ROD and FRED-OX, where \texttt{T\_coolant} comes from a user-specified time series or an external thermal-hydraulics code.
  \item[Robin (\texttt{isCladTempOuterBoundaryPrescribed = false}).] An instantaneous energy balance is imposed: the conductive flux arriving at the outer cladding ring equals the convective flux leaving to the coolant, $Q_{n_c-2\to n_c-1} = h_{cool}\cdot 2\pi r_{co}\cdot(T_{n_c-1} - T_{cool})$, where \texttt{h\_cool} [W/(m$^2$\,K)] is the cladding-to-coolant heat-transfer coefficient.  The Robin mode is used by FRED-M-Na when the in-built sub-channel model is active (Section~\ref{sec:subchannel_plat}). Application residual classes should not set this flag directly; instead, they should use the \texttt{SubchannelMode} abstraction described below.
\end{description}

\paragraph{StressStrain.}
Implements the thermo-elastic stress-strain equations in cylindrical geometry, including fuel swelling (via \texttt{AxialLayerState::efs}) and cladding hoop creep (via \texttt{AxialLayerState::ec}).  The equation count is \coderef{stress_strain_neq}: $6 n_f + 6 n_c + 2$ equations.  The method \coderef{stress_strain_residual} fills the mechanics sub-vector.

The stress equations are entirely algebraic: all strain and stress fields appear as AE variables in the IDA system.  Only $n_f$ fuel temperatures and $n_c - 1$ cladding temperatures are differential.

\subsection{Gap Utilities}
\label{sec:gap_utils}

\paragraph{GapStateManager.}
\coderef{gap_state_manager} tracks the open/closed state and the axial strain offset for each axial layer.  When the gap closes, the constraint $\varepsilon_z^\text{fuel} - \varepsilon_z^\text{clad} = C$ (where $C$ is frozen at the moment of contact) is imposed algebraically.  This offset is cleared on reopening.  All residual assemblers share a single instance of \texttt{GapStateManager} owned by \texttt{RodResiduals}.

\paragraph{GapPressureModel.}
\coderef{gap_pressure_model} formulates the global gas pressure algebraic equation using the gas bond model: ideal gas law using the fill-gas inventory $\mu_0$ and the temperature-weighted free volume $V/T$ summed over the plenum, annular gap, and central void.
FRED-OX does not use \texttt{GapPressureModel}; it incorporates the gas pressure equation directly into its global residual block alongside the fission gas ODEs.

\subsection{Sub-Channel Coolant Interface}
\label{sec:subchannel_plat}

\paragraph{SubchannelCoolantProperties.}
An abstract interface (\texttt{platform/SubchannelCoolantProperties.hpp}) that applications must implement to provide sodium (or other coolant) thermophysical properties: specific heat \texttt{cp(T)}, density \texttt{rho(T)}, thermal conductivity \texttt{k(T)}, and dynamic viscosity \texttt{mu(T)}.

\paragraph{SubchannelMode.}
The abstract class \texttt{SubchannelMode} (\texttt{platform/SubchannelMode.hpp}) defines the interface for computing per-layer coolant temperatures and cladding-to-coolant HTCs from inlet conditions.  The single pure virtual method \texttt{updateCoolantField(states, t)} is called once per residual evaluation, before the per-layer loop in \texttt{RodResiduals}.  After the call:
\begin{itemize}
  \item \texttt{subchannelMode.T\_co(j)} returns the bulk coolant temperature at layer~$j$ [K].
  \item \texttt{subchannelMode.htc(j)} returns the cladding-to-coolant HTC at layer~$j$ [W/(m$^2$\,K)].
\end{itemize}
When a \texttt{SubchannelMode} is registered via \texttt{RodResiduals::setSubchannelMode(\&mode)}, the thermal residual call for each layer automatically switches to the Robin BC mode: \texttt{isCladTempOuterBoundaryPrescribed = false} with the HTC from \texttt{htc(j)} and the coolant temperature from \texttt{T\_co(j)}.

\paragraph{Implementing a sub-channel model.}
Concrete subclasses inherit \texttt{SubchannelMode} and override \texttt{updateCoolantField}.  The FRED-M-Na sodium sub-channel model (\texttt{FredMNaSubchannelMode}) serves as the reference implementation: it performs an axial energy balance using the one-dimensional steady-state equation and evaluates the cladding-to-coolant HTC from Péc\-let-number correlations (Mikityuk or Subbotin; see the Theory Manual, Section~4 for the equations). Only FRED-M-Na currently uses the sub-channel mode; FRED-ROD and FRED-OX continue to use the Dirichlet BC with a prescribed coolant temperature.

\subsection{RodResiduals — the NVI Residual Orchestrator}
\label{sec:rod_residuals}

\texttt{RodResiduals} (\coderef{rod_residuals_class}) is the central abstract base class for DAE residual assembly.  It owns \texttt{HeatConduction}, \texttt{StressStrain}, and \texttt{GapStateManager}, and drives a structured assembly loop.  No application code should call the physics blocks directly.

\subsubsection*{The assembly loop}

The public entry point \coderef{rod_residuals_nvi} is a Non-Virtual Interface (NVI): it calls a fixed sequence of virtual hooks, allowing applications to inject physics at precisely defined points without re-implementing the shared loop logic.  The call sequence for each IDA residual evaluation is:

\begin{enumerate}
  \item \textbf{\texttt{unpackState(y, yp)}} — copy the flat IDA vector into the application's typed state struct.
  \item \textbf{\texttt{computeGlobalResiduals(t, yp, r)}} — fill the \texttt{globalOffset()} equations at the front of \texttt{r} (pressure AE, fission gas ODEs, \ldots).
  \item For each axial layer $j = 0 \ldots n_z - 1$:
        \begin{enumerate}
          \item \textbf{\texttt{prepareLayer(j)}} — optional per-layer pre-work (default: no-op; FRED-OX uses this to propagate gap state from \texttt{GapStateManager} into the typed layer state).
          \item \textbf{\texttt{computeThermalResiduals(j, \ldots)}} or \textbf{\texttt{freezeThermalResiduals(j, \ldots)}} depending on the \texttt{m\_heat\_on} toggle.
          \item \textbf{\texttt{computeIrradiationResiduals(j, \ldots)}} — burnup and swelling ODEs (default: no-op for FRED-ROD).
          \item \texttt{StressStrain::computeResiduals} via the base class.
        \end{enumerate}
\end{enumerate}

\subsubsection*{Equation counts}

The per-layer equation layout is controlled by three members defined at \coderef{rod_residuals_neq_layout}:
\begin{center}
\begin{tabular}{lll}
  \toprule
  Member & Equations & Value \\
  \midrule
  \texttt{m\_neq\_th}   & thermal block   & $4 + n_f + n_c$ \\
  \texttt{m\_neq\_irr}  & irradiation block & $0$ (ROD), $1 + n_f + n_c$ (OX) \\
  \texttt{m\_neq\_mech} & mechanics block & $6 n_f + 6 n_c + 2$ \\
  \texttt{m\_neq\_j}    & total per layer & $m\_neq\_th + m\_neq\_irr + m\_neq\_mech$ \\
  \bottomrule
\end{tabular}
\end{center}

\subsubsection*{Virtual hooks summary}

The hooks declared at \coderef{rod_residuals_hooks} that a subclass must or may override are listed in Table~\ref{tab:hooks}.

\begin{table}[h]
\centering
\caption{Virtual hooks in \texttt{RodResiduals}.}
\label{tab:hooks}
\small
\begin{tabular}{llp{6cm}}
  \toprule
  Hook & Required & Purpose \\
  \midrule
  \texttt{unpackState}               & yes & flat IDA $\mathbf{y}/\dot{\mathbf{y}}$ → typed state \\
  \texttt{globalOffset}              & yes & number of global equations \\
  \texttt{globalPressure}            & yes & current gas pressure (for mechanics BC) \\
  \texttt{layerState}                & yes & per-layer state reference \\
  \texttt{computeGlobalResiduals}    & yes & fill global equation block \\
  \texttt{computeThermalResiduals}   & yes & fill thermal block for layer $j$ \\
  \texttt{onGapClosed}               & yes & record closure offset in typed state \\
  \texttt{freezeThermalResiduals}    & no  & override when heat is off \\
  \texttt{computeIrradiationResiduals} & no & fill irradiation block (default: no-op) \\
  \texttt{prepareLayer}              & no  & per-layer pre-work (default: no-op) \\
  \bottomrule
\end{tabular}
\end{table}

\subsubsection*{Solution-vector layout}
\label{subsec:yvec_layout}

The residual function is implemented in \texttt{RodResiduals} (\texttt{src/platform/RodResiduals.cpp}) following the NVI pattern described above.  The solution vector is organised in three blocks:

\begin{equation}
  \mathbf{y} = \Bigl[
    \underbrace{\mathbf{y}_\mathrm{glob}}_{\text{global}}
    \;\Big|\;
    \underbrace{\mathbf{y}_{j=1}}_{\text{layer }1}
    \;\Big|\;
    \cdots
    \;\Big|\;
    \underbrace{\mathbf{y}_{j=n_z}}_{\text{layer }n_z}
  \Bigr]^\top,
  \qquad
  N = N_\mathrm{glob} + n_z \cdot N_j.
  \label{eq:y_layout}
\end{equation}

\paragraph{Global block} ($N_\mathrm{glob}$ entries, returned by \texttt{neqGlobal()}).
Application-defined; for FRED-OX and FRED-M-Na this is three entries (fission gas generated $n_{gen}$, fission gas released $n_{rel}$, gas pressure $p_{gas}$); for FRED-ROD there is gas pressure only.

\paragraph{Per-layer block} ($N_j$ entries per axial layer~$j$).
The per-layer block is subdivided into three sub-blocks:

\begin{enumerate}
  \item \textbf{Thermal sub-block} ($N_\mathrm{th} = 4 + n_f + n_c$ entries):
    \begin{table}[h]
      \centering
      \begin{tabular}{lll}
        \toprule
        Symbol & Description & Type \\
        \midrule
        $q_{v,j}^{\prime\prime\prime}$ & Volumetric power density & AE (table-lookup BC) \\
        $\delta_j$            & Fuel--cladding gap width         & AE \\
        $p_{\mathrm{fc},j}$  & Fuel--cladding contact pressure   & AE \\
        $h_{\mathrm{gap},j}$ & Gap conductance                   & AE \\
        $T_{f,j,i}$, $i=1\ldots n_f{-}1$ & Fuel ring temperatures & ODE \\
        $T_{c,j,i}$, $i=1\ldots n_c{-}1$ & Inner cladding temperatures & ODE \\
        $T_{\mathrm{co},j}$  & Outer cladding surface temperature & AE (BC) \\
        \bottomrule
      \end{tabular}
      \caption{Thermal sub-block variables for axial layer~$j$.}
      \label{tab:thermal_block}
    \end{table}

  \item \textbf{Irradiation sub-block} ($N_\mathrm{irr}$ entries, application-defined via the \texttt{neq\_irr} constructor argument). For FRED-OX: burnup~$B_j$ (ODE), per-node fuel swelling strains~$\varepsilon_{s,i}$ (ODE), cladding creep strains~$\varepsilon_{c,i}$ (ODE). For FRED-ROD: empty.

  \item \textbf{Mechanics sub-block} ($N_\mathrm{mech}$ entries, fixed by \texttt{StressStrain}). Fuel and cladding thermal, elastic, and total strains; hoop, radial, and axial stresses. All are algebraic (AE).
\end{enumerate}

\paragraph{Variable-type (\texttt{id}) array.}
Each concrete \texttt{RodResiduals} subclass overrides \texttt{fillIdArray} to set $\mathtt{id}_k = 1$ for differential variables and $\mathtt{id}_k = 0$ for algebraic variables, following the sub-block structure above.  This array is passed to IDA via \texttt{IDASetId}, which uses it to: (i) suppress algebraic variables from the LTE error test, and (ii) partition $\mathbf{y}$ for \texttt{IDACalcIC}.

\subsection{FredSolverBase — the IDA Driver}
\label{sec:solver_base}

\texttt{FredSolverBase} (\coderef{fred_solver_base}) provides the shared SUNDIALS machinery used by all solver drivers.  It eliminates duplication of SUNDIALS context management, N\_Vector allocation, output storage, and gap-event handling.

\subsubsection*{Subclass setup hooks}

Five pure virtual hooks (declared at \coderef{solver_setup_hooks}) must be implemented by every solver subclass.  They are called once by \coderef{solver_setup_sundials}:

\begin{center}
\begin{tabular}{ll}
  \toprule
  Hook & Returns \\
  \midrule
  \texttt{getSolverNeq()}      & total number of DAE equations \\
  \texttt{getSolverResFn()}    & \texttt{IDAResFn} callback (usually \texttt{RodResiduals::idaResidual}) \\
  \texttt{getSolverUserData()} & user-data pointer passed to IDA \\
  \texttt{fillSolverIdArray(id, neq)} & mark each entry as ODE (1) or AE (0) \\
  \texttt{packSolverState(y)}  & pack the initial application state into the IDA $\mathbf{y}$ vector \\
  \bottomrule
\end{tabular}
\end{center}

\subsubsection*{setupSundials initialisation sequence}

The full initialisation sequence performed by \coderef{solver_setup_sundials} is:
\begin{enumerate}
  \item Create a SUNDIALS context: \texttt{SUNContext\_Create}.
  \item Allocate serial \texttt{N\_Vector} objects for $\mathbf{y}$, $\dot{\mathbf{y}}$, and the tolerance vector.
  \item Create IDA memory: \texttt{IDACreate}.
  \item Register the residual callback: \texttt{IDAInit(m\_ida\_mem, RodResiduals::idaResidual, \ldots)}.
  \item Set scalar relative and vector absolute tolerances: \texttt{IDASVtolerances}.
  \item Create and attach a dense $N \times N$ matrix (\texttt{SUNDenseMatrix}) and dense direct linear solver (\texttt{SUNLinSol\_Dense}).
  \item Create and attach the Newton nonlinear solver (\texttt{SUNNonlinSol\_Newton}).
  \item Set the maximum number of internal steps: \texttt{IDASetMaxNumSteps}.
  \item Register gap root functions: \texttt{IDARootInit} with \texttt{RodResiduals::gapRoot} as the root function callback (one root function per axial layer).
\end{enumerate}

\subsubsection*{Initial condition calculation (calcIC)}

In \texttt{FredSolverBase::calcIC} (\texttt{src/platform/FredSolverBase.cpp}), the call sequence is:
\begin{enumerate}
  \item Pack the initial physical state into the IDA y-vector via \texttt{packSolverState}.
  \item Set the id array via \texttt{IDASetId}.
  \item Suppress algebraic variables from the LTE test via \texttt{IDASetSuppressAlg}.
  \item Call \texttt{IDACalcIC(m\_ida\_mem, IDA\_YA\_YDP\_INIT, t\_ic)}, where $t_\mathrm{ic}$ is the first output time.
  \item Retrieve the consistent solution via \texttt{IDAGetConsistentIC}.
\end{enumerate}
Because the IC calculation is itself a nonlinear solve, FRED 2.0 relaxes the Newton tolerances for this step (increasing \texttt{IDASetMaxNumStepsIC} and \texttt{IDASetMaxNumItersIC}, and loosening \texttt{IDASetNonlinConvCoefIC}) to improve robustness for problems with a large initial thermal gradient.

\subsubsection*{The time loop}

\coderef{solver_run_time_loop} advances IDA in \texttt{IDA\_ONE\_STEP} mode, calling
\begin{equation*}
  \texttt{IDASolve(m\_ida\_mem,\; t\_far,\; \&t\_ret,\; y,\; yp,\; IDA\_ONE\_STEP)}
\end{equation*}
repeatedly until $t_\mathrm{ret} \geq t_\mathrm{end}$.  At each accepted step the following virtual hooks are called:
\begin{itemize}
  \item \textbf{\texttt{unpackCurrentState()}} — copy IDA $\mathbf{y}/\dot{\mathbf{y}}$ into the solver's typed state.
  \item \textbf{\texttt{storeOutput(t)}} — append the current state to output vectors (app-specific).
  \item \textbf{\texttt{logStepOutput(t)}} — console progress line (optional override).
\end{itemize}
After each call the loop checks whether a root was detected and dispatches the appropriate gap-event handler, then stores output at the scheduled output times.

\subsubsection*{Gap event handling}

When \texttt{IDASolve} returns \texttt{IDA\_ROOT\_RETURN}, the time loop dispatches:
\begin{itemize}
  \item \textbf{Closure} (\texttt{rootsfound[j] = -1}): calls \texttt{handleGapClosed(j)}, which records the axial offset at closure in \texttt{GapStateManager::applyGapClosed} and switches the residual branch to the closed-gap equations.
  \item \textbf{Reopening} (\texttt{rootsfound[j] = +1}): calls \texttt{handleGapReopened(j)}, which restores the open-gap residual branch.
\end{itemize}

\paragraph{Why \texttt{gapOpen} must not be modified inside the residual.}
IDA calls the residual function an unbounded number of times per step with speculative, possibly-rejected trial values.  The contact-state flag \texttt{gapOpen[j]}, managed by \texttt{GapStateManager}, must not be modified during any of these residual evaluations.  Only the event handler in the main time loop, executing after \texttt{IDASolve} returns \texttt{IDA\_ROOT\_RETURN}, may change the contact state.  Violation of this rule would corrupt the root-function sign bracket that IDA's bisection algorithm relies on, leading to missed or spurious events.

\paragraph{Structural restart after a contact event.}
A contact-state transition changes the algebraic regime of the DAE: the closed-gap branch introduces a different contact-pressure constraint and pins $\delta_j = \delta_\mathrm{contact}$, while the open-gap branch releases this pin and sets $p_\mathrm{fc,j} = 0$.  Continuing with the multistep BDF history built for the pre-event branch can cause Newton failures or degrade accuracy.  FRED 2.0 therefore applies a structural restart at every contact event, implemented in \texttt{FredSolverBase::afterGapEventsHandled} $\to$ \texttt{reinitAfterEvent} (\texttt{src/platform/FredSolverBase.cpp}):
\begin{enumerate}
  \item Switch the contact state in \texttt{GapStateManager}.
  \item Call \texttt{IDAReInit} at the event time to clear the BDF history.
  \item Reassert the id array via \texttt{IDASetId}.
  \item Recompute consistent initial conditions via \texttt{IDACalcIC}.
  \item Retrieve the consistent state via \texttt{IDAGetConsistentIC} and resume integration from the event time.
\end{enumerate}

% ===========================================================================
\section{Python Bindings Framework}
\label{sec:python_bindings}
% ===========================================================================

FRED~2.0 exposes all solver-facing interfaces to Python via pybind11.  The binding layer serves two distinct purposes: input validation and the experimental correlation interface.

\subsection{User-facing principle: Python is the entry point}

Users should \emph{never} construct or manipulate C++ objects directly.  All inputs --- geometry, material correlations, boundary conditions, solver settings --- flow through the Python-bound objects.  This boundary provides a single, auditable entry point for all simulations and separates user-facing configuration from internal solver mechanics.

\subsection{Experimental correlation interface}

The abstract material interfaces (\texttt{FuelPelletMaterial}, \texttt{CladdingMaterial}, \texttt{GapMaterial}) are bound with pybind11 \emph{trampolines}.  A trampoline is a thin C++ shim that, instead of dispatching to a C++ virtual method, calls back into the Python runtime.  This allows a pure-Python subclass to override any material method: when the IDA solver calls \texttt{fuel.thermalConductivity(T)} on a Python-defined material object, the Python implementation runs with no C++ recompilation.

\paragraph{Intended use cases.}
\begin{description}
  \item[Rapid prototyping.]  Implement a new $k(T)$ or $\varepsilon_\mathrm{th}(T)$ correlation in a Python class.  Run the full solver, inspect peak temperature and stress outputs, and iterate.  The cycle is minutes rather than the hours a C++ rebuild requires.  This is the \textbf{recommended first step} for any new correlation.
  \item[Sensitivity studies.]  Accept free parameters in the class constructor; sweep them with a Python loop or \texttt{numpy.linspace} without touching any input file.
  \item[Model calibration / optimisation.]  Connect solver outputs (\texttt{peak\_fuel\_temperature()}, \texttt{gap\_width()}, etc.) to a \texttt{scipy.optimize} routine; let it tune correlation coefficients automatically.
  \item[Irradiation-adjusted properties.]  Inject burnup- or dose-dependent corrections that read from a pre-computed lookup table, with no changes to the C++ solver state.
\end{description}

Once numerically validated, a Python correlation can be ported to C++ and registered as a built-in material for production use (Section~\ref{sec:fred_rod_materials}).

\subsection{Binding architecture}

The top-level module is assembled in \texttt{src/bindings.cpp}, which calls one \texttt{bind\_*} function per application.  Each function lives in the application's own \texttt{bindings.cpp}:

\begin{center}
\begin{tabular}{ll}
  \toprule
  Function & File \\
  \midrule
  \texttt{bind\_fred\_rod(m)}    & \texttt{src/apps/fred\_rod/bindings.cpp} \\
  \texttt{bind\_fred\_ox(m)}     & \texttt{src/apps/fred\_ox/bindings.cpp} \\
  \texttt{bind\_fred\_m\_na(m)}  & \texttt{src/apps/fred\_m\_na/bindings.cpp} \\
  \bottomrule
\end{tabular}
\end{center}

\texttt{bind\_fred\_rod} must be called first because it registers the platform abstract types (\texttt{FuelPelletMaterial}, \texttt{CladdingMaterial}, \texttt{GapMaterial}, \texttt{FuelRodGeometry}) with their trampolines.  The OX and M-Na bindings then register concrete subclasses of those types; pybind11 requires that the base class is already registered.

\paragraph{Keep-alive semantics.}
Solvers hold \emph{references} to material objects, not copies.  The \texttt{py::keep\_alive<1,N>()} annotations in each solver constructor inform the garbage collector that the material must remain alive as long as the solver does.  Binding authors must apply \texttt{keep\_alive} to every reference argument of any solver constructor.

% ===========================================================================
\section{Application: FRED-ROD}
\label{sec:fred_rod}
% ===========================================================================

FRED-ROD is the baseline application.  It solves thermal conduction and thermo-elastic stress-strain without any irradiation physics.  Users are expected to supply their own \texttt{FuelPelletMaterial} and \texttt{CladdingMaterial} implementations for experimental workflows.

\subsection{FredRodResiduals}
\label{sec:rod_res}

\coderef{fred_rod_residuals} inherits \texttt{RodResiduals} and implements the minimum set of required hooks.  The y-vector layout is documented at \coderef{fred_rod_yvec_layout}:

\begin{center}
\begin{tabular}{ll}
  \toprule
  Index range & Variable \\
  \midrule
  \texttt{y[0]} & \texttt{gpres} — global gas pressure AE \\
  \texttt{y[1 + j*neq\_j \ldots 1+(j+1)*neq\_j-1]} & layer $j$ block (thermal + mechanics) \\
  \bottomrule
\end{tabular}
\end{center}

Within each layer, the thermal sub-block (\texttt{m\_neq\_th} entries) comes first, followed by the mechanics sub-block (\texttt{m\_neq\_mech} entries). There is no irradiation sub-block (\texttt{m\_neq\_irr} = 0).

The overridden hook \coderef{fred_rod_thermal_hook} delegates the global pressure equation to \texttt{GapPressureModel::residual}, and the thermal equations to \texttt{HeatConduction::computeResiduals}.

\subsection{FredRodSolver}
\label{sec:rod_solver}

\coderef{fred_rod_solver} inherits \texttt{FredSolverBase}.  The main entry point is \coderef{fred_rod_run}.  Internally it:

\begin{enumerate}
  \item Calls \texttt{initState()} — allocates \texttt{RodState}, computes initial gap widths, calls \texttt{FredRodResiduals::initAlgebraicState()} to set AE initial values, then \texttt{solveMechanicalIC()} for a Newton pre-solve on the mechanical equilibrium.
  \item Calls \texttt{setupSundials()} (fills IDA id array, packs state into $\mathbf{y}$, sets tolerances).
  \item Calls \texttt{setupGapRoots(nz, RodResiduals::gapRoot)} to register gap event detection.
  \item Calls \texttt{calcIC(t\_start)} for consistent initial conditions.
  \item Calls \texttt{runTimeLoop(tend, dtout, all\_steps)}.
\end{enumerate}

After a gap event, the FRED-ROD override of \texttt{afterGapEventsHandled()} additionally re-solves the mechanical equilibrium before calling \texttt{reinitAfterEvent()}.

\paragraph{Result accessors.}
The solver exposes per-step vectors: \texttt{gapWidth()}, \texttt{contactPressure()}, \texttt{cladOuterHoopStress()}, \texttt{fuelOuterRadius()}, \texttt{cladInnerRadius()}, and the shared \texttt{peakFuelTemperature()} from the base class.

% ===========================================================================
\subsection{Material Implementation: FRED-ROD as Reference}
\label{sec:fred_rod_materials}
% ===========================================================================

FRED-ROD is the canonical application for adding new material correlations. The recommended workflow has two stages: a Python prototype (no compilation required) followed by an optional C++ port for production performance.

% -----------------------------------------------------------------------
\subsubsection{Stage 1 --- Python material (recommended first step)}
% -----------------------------------------------------------------------

Subclass \texttt{fred\_rod.FuelPelletMaterial}, \texttt{fred\_rod.CladdingMaterial}, or \texttt{fred\_rod.GapMaterial} in Python and override every abstract method. The instance is passed directly to \texttt{FredRodSolver}; no C++ changes are needed.

A complete working example is \texttt{examples/fred\_rod/python\_materials/run.py}. The condensed pattern for a fuel pellet is:

\begin{lstlisting}[language=Python]
import fred_rod as fred

class MyFuel(fred.FuelPelletMaterial):
    def thermalConductivity(self, T):        # [W/(m·K)]
        return 1.0 / (0.045 + 2.8e-4 * T)
    def heatCapacity(self, T):               # [J/(kg·K)]
        return 260.0 + 0.10 * T
    def thermalExpansionStrain(self, T):     # [-]
        return 1.0e-5 * (T - 293.15)
    def youngsModulus(self, T, density):     # [MPa]
        return 2.0e5
    def poissonRatio(self):        return 0.30
    def referenceDensity(self):    return 10500.0
    def theoreticalDensity(self):  return 10960.0
    def meltingTemperature(self):  return 3100.0

geom = fred.FuelRodGeometry()
geom.nf = 4;  geom.nc = 3;  geom.nz = 1
geom.rfi0 = 0.0;  geom.rfo0 = 4.2e-3
geom.rci0 = 4.3e-3;  geom.rco0 = 4.9e-3
geom.dz0 = [0.10];  geom.vgp = 1e-6
geom.ruff = 5e-6;  geom.rufc = 5e-6
geom.build()

solver = fred.FredRodSolver(geom, MyFuel(), fred.AIM1(), fred.He())
solver.set_power_density_history([0, 1000], [1.5e8, 1.5e8])
solver.set_coolant_temperature([0, 1000], [620.0, 620.0])
solver.set_initial_temperature(620.0)
solver.set_initial_gas_pressure(0.1)
solver.set_tolerances(1e-6, 1e-8)
solver.run(1000.0, 100.0)
print(solver.peak_fuel_temperature())
\end{lstlisting}

For a liquid-bond \texttt{GapMaterial} the pattern is:

\begin{lstlisting}[language=Python]
class MyLiquidBond(fred.GapMaterial):
    def gapConductivity(self, T):   # k [W/(m·K)] of the bond liquid
        return 110.0 - 0.065 * T
    def is_gas_bond(self):
        return False
    def clamp_conductance(self, h, gap_closed):
        return max(h, 1.0e6) if gap_closed else h  # enforce 1 MW/m2K floor
\end{lstlisting}

\paragraph{Numerical validation checklist.}
\begin{enumerate}
  \item Plot $k(T)$, $c_p(T)$, $\varepsilon_\mathrm{th}(T)$ over the full temperature range before running the solver.
  \item Run with \texttt{set\_enable\_stress\_strain(False)} first to isolate the thermal solution.
  \item Tighten tolerances to $\mathtt{rtol}=10^{-8}$, $\mathtt{atol}=10^{-10}$ and confirm that peak fuel temperature changes by less than 0.1~K --- this verifies the correlation is smooth enough for the IDA numerical Jacobian.
  \item Enable stress-strain and verify that gap width, contact pressure, and cladding hoop stress evolve plausibly.
\end{enumerate}

% -----------------------------------------------------------------------
\subsubsection{Stage 2 --- C++ material}
% -----------------------------------------------------------------------

Once the Python prototype is validated, port it to C++ for production use.

\paragraph{(a) Implement the class.}
Create \texttt{src/apps/fred\_rod/fuelpelletmaterial/MyFuel.hpp} and \texttt{MyFuel.cpp}.  Inherit from \texttt{fred::FuelPelletMaterial} and implement all pure virtual methods:

\begin{lstlisting}
// MyFuel.hpp
#pragma once
#include "platform/FuelPelletMaterial.hpp"
namespace fred {
class MyFuel : public FuelPelletMaterial {
public:
    double thermalConductivity(double T) const override;
    double heatCapacity(double T) const override;
    double thermalExpansionStrain(double T) const override;
    double youngsModulus(double T, double density) const override;
    double poissonRatio() const override;
    double referenceDensity() const override;
    double theoreticalDensity() const override;
    double meltingTemperature() const override;
};
} // namespace fred
\end{lstlisting}

\begin{lstlisting}
// MyFuel.cpp
#include "MyFuel.hpp"
namespace fred {
double MyFuel::thermalConductivity(double T) const {
    return 1.0 / (0.045 + 2.8e-4 * T);
}
// ... remaining methods mirror the Python prototype
} // namespace fred
\end{lstlisting}

Add \texttt{MyFuel.cpp} to the \texttt{MAT\_SRCS} variable in the \texttt{Makefile}.

\paragraph{(b) Register Python bindings.}
In \texttt{src/apps/fred\_rod/bindings.cpp}, include the header and add the class registration inside \texttt{bind\_fred\_rod()}:

\begin{lstlisting}
#include "apps/fred_rod/fuelpelletmaterial/MyFuel.hpp"
// inside bind_fred_rod(py::module_& m):
py::class_<MyFuel, FuelPelletMaterial>(m, "MyFuel")
    .def(py::init<>(), "My custom fuel pellet.")
    .def("thermal_conductivity",
         &MyFuel::thermalConductivity,    py::arg("T"))
    .def("heat_capacity",
         &MyFuel::heatCapacity,           py::arg("T"))
    .def("thermal_expansion_strain",
         &MyFuel::thermalExpansionStrain, py::arg("T"))
    .def("youngs_modulus",
         &MyFuel::youngsModulus,          py::arg("T"), py::arg("density"))
    .def("poisson_ratio",          &MyFuel::poissonRatio)
    .def("reference_density",      &MyFuel::referenceDensity)
    .def("theoretical_density",    &MyFuel::theoreticalDensity)
    .def("melting_temperature",    &MyFuel::meltingTemperature);
\end{lstlisting}

Rebuild with \texttt{make pymodule}.  The Python call site (\texttt{fred.MyFuel()}) is identical to the prototype --- simulation scripts need no changes.

% ===========================================================================
\section{Application: FRED-OX}
\label{sec:fred_ox}
% ===========================================================================

FRED-OX augments FRED-ROD with MOX irradiation physics: burnup, fuel swelling, fission gas generation and release, and irradiation creep of the cladding.  It uses the same \texttt{HeatConduction} and \texttt{StressStrain} blocks from the platform layer, wrapped by its own residual assembler.

\subsection{Extended State}
\label{sec:ox_state}

\coderef{fred_ox_layer_state} (\texttt{FredOxLayerState}) extends \texttt{AxialLayerState} with three additional fields:
\begin{itemize}
  \item \texttt{bup} — local burnup [MWd/kgU].
  \item \texttt{defs[nf]} — per-node fuel swelling rate [1/s].
  \item \texttt{dec[nc]} — per-ring cladding hoop creep rate [1/s].
\end{itemize}
The base fields \texttt{efs[nf]} (fuel swelling strain) and \texttt{ec[nc]} (cladding creep strain) are declared in \texttt{AxialLayerState} so that \texttt{StressStrain} can consume them generically.

\coderef{fred_ox_rod_state} (\texttt{FredOxRodState}) aggregates \texttt{FredOxLayerState} layers and global gas-phase quantities: \texttt{gpres}, \texttt{fggen}, \texttt{fgrel}, and their rates.

\subsection{FredOxResiduals}
\label{sec:ox_res}

\coderef{fred_ox_residuals} adds three global equations (\coderef{fred_ox_global_residuals}):
\begin{itemize}
  \item \texttt{fggen} ODE — total fission gas generated [mol/s].
  \item \texttt{fgrel} ODE — total fission gas released [mol/s].
  \item \texttt{gpres} AE — ideal gas law with fission gas: $p = (\mu_0 + n_\text{rel}) R / (V/T) \times 10^{-6}$ [MPa].
\end{itemize}

Per layer, the irradiation sub-block (\coderef{fred_ox_irradiation_hook}) adds $1 + n_f + n_c$ equations:
\begin{itemize}
  \item One burnup ODE: $\dot{b}_y = \texttt{qqv} / \rho_0$ where $b_y = b \times 8.64 \times 10^{10}$ [J/kgU].
  \item $n_f$ fuel swelling ODEs: $\dot{\varepsilon}^s_f[i] = f_\text{swm} \cdot \dot{\varepsilon}^s_\text{MATPRO}(T_f[i], \dot{b})$.
  \item $n_c$ cladding creep ODEs: $\dot{\varepsilon}^c[i] = \texttt{CladdingMaterial::creepRate}(T_c[i], \sigma_h[i])$.
\end{itemize}

\paragraph{FredOxGapMaterial.}
The FRED-OX gap material accounts for He/Xe/Kr gas mixture conductivity (from fission gas release), the Lanning--Hann temperature-jump model, and radiation conductance with KIT emissivities.  The fission gas composition is fixed at 88.46\%~Xe, 7.69\%~Kr, 3.85\%~He by mole.  The solver updates the gap material state setters (\texttt{setFissionGasRelease()}, \texttt{setGasPressure()}, \texttt{setBurnup()}) at each output step before evaluating gap conductance.

\subsection{FredOxSolver}
\label{sec:ox_solver}

\coderef{fred_ox_solver} mirrors the structure of \texttt{FredRodSolver} but omits the mechanical pre-solve step (the gap-event path in FRED-OX does not re-solve mechanical IC separately).  It also overrides \texttt{logStepOutput()} to print gas pressure and fission gas release alongside the time.

Result accessors: \texttt{gasPressure()}, \texttt{fgGenerated()}, \texttt{fgReleased()}, \texttt{gapWidth()}, \texttt{burnup()}.

% ===========================================================================
\section{How to Add a New Material}
\label{sec:new_material}
% ===========================================================================

\subsection{New fuel pellet material for FRED-ROD}

\begin{enumerate}
  \item Create \texttt{src/apps/fred\_rod/fuelpelletmaterial/MyFuel.hpp} and \texttt{MyFuel.cpp}.
  \item Inherit from \texttt{fred::FuelPelletMaterial} (\coderef{fuel_pellet_iface}) and implement all pure virtual methods.
  \item Add \texttt{MyFuel.cpp} to the \texttt{MAT\_SRCS} list in \texttt{Makefile}.
  \item Expose the class to Python by adding a \texttt{py::class\_<MyFuel>} binding in \texttt{src/apps/fred\_rod/bindings.cpp}.
\end{enumerate}

The conductivity, heat capacity, thermal expansion, and elastic moduli must be self-consistent with the units used throughout: W/(m$\cdot$K), J/(kg$\cdot$K), dimensionless strain, and MPa respectively.

\subsection{New cladding material with irradiation creep}

Follow the same steps, but inherit from \texttt{fred::CladdingMaterial} (\coderef{cladding_iface}) and override the optional \coderef{cladding_creep_hook}.  If the material is intended for FRED-OX, place it under \texttt{src/apps/fred\_ox/claddingmaterial/} and include it in the OX bindings.  The irradiation residual block (\coderef{fred_ox_irradiation_hook}) will automatically call \texttt{creepRate(T, sigh)} when assembling the cladding creep ODE.

\subsection{Liquid-bond gap material}

Inherit from \texttt{GapMaterial} (\coderef{gap_material_iface}) and:
\begin{enumerate}
  \item Implement \texttt{gapConductivity(T)} to return the liquid's thermal conductivity $k$ [W/(m$\cdot$K)] at average gap temperature $T$ [K].
  \item Override \coderef{gap_bond_selector} to return \texttt{false}.
  \item Optionally override \texttt{clampConductance(h, gap\_closed)} to enforce minimum/maximum conductance bounds.
\end{enumerate}
The platform will then compute $h = k(T) / \max(\delta, \sqrt{r_\mathrm{uf}^2+r_\mathrm{uc}^2})$ and pass the result through \texttt{clampConductance} before using it. There is no method \texttt{gapConductance(T)} in the interface; the platform always derives $h$ from $k$.

% ===========================================================================
\section{How to Add a New Application}
\label{sec:new_app}
% ===========================================================================

A new application (e.g.\ FRED-M-Na for metallic fuel) needs four files:

\begin{enumerate}
  \item \textbf{State struct} — an optional struct inheriting \texttt{AxialLayerState} (\coderef{axial_layer_state}) and/or \texttt{RodState} (\coderef{rod_state}) that adds irradiation-specific fields.
  \item \textbf{Residuals class} — inherits \texttt{RodResiduals} (\coderef{rod_residuals_class}) and implements the mandatory hooks from Table~\ref{tab:hooks}.  The constructor must pass \texttt{neq\_irr} to the protected \texttt{RodResiduals} constructor to declare the size of the irradiation block.
  \item \textbf{Solver class} — inherits \texttt{FredSolverBase} (\coderef{fred_solver_base}), implements the five setup hooks from \coderef{solver_setup_hooks}, and provides a \texttt{run()} method that calls \texttt{setupSundials()}, \texttt{calcIC()}, and \texttt{runTimeLoop()}.
  \item \textbf{Bindings} — a \texttt{bind\_<appname>} function registered in \texttt{src/bindings.cpp}.
\end{enumerate}

\paragraph{Key invariant.}
The irradiation block occupies y-vector indices $[\texttt{globalOffset()} + j \cdot \texttt{m\_neq\_j} + \texttt{m\_neq\_th}$ to $\texttt{globalOffset()} + j \cdot \texttt{m\_neq\_j} + \texttt{m\_neq\_th} + \texttt{m\_neq\_irr} - 1]$ for layer~$j$.  The mechanics block always follows.  Packing and unpacking must be consistent with this layout.

\paragraph{Platform files are not to be modified.}
All shared physics is in \texttt{platform/} and must not be altered for application-specific needs.  If a physics block genuinely needs extension, discuss with the platform maintainer.

\paragraph{Worked example: extending mechanics without touching the platform (FRED-M-Na).}
\texttt{StressStrain::computeResiduals} (\coderef{stress_strain_residual}) is
non-virtual, and \texttt{RodResiduals::computeResiduals}
(\coderef{rod_residuals_nvi}) — the shared per-layer assembly loop that calls
it — is also non-virtual, so neither can be overridden by a subclass without
editing the platform header. FRED-M-Na needed a three-state (\texttt{open}/
\texttt{soft}/\texttt{clos}) mechanics model for U-Pu-Zr metallic fuel instead
of the shared two-state (\texttt{open}/\texttt{closed}) model, without
affecting FRED-ROD/FRED-OX. The pattern used, in case a future application
needs the same kind of extension:
\begin{enumerate}
  \item Fork the mechanics class as a \texttt{.cpp}-only copy —
    \texttt{FredMNaStressStrain} (\coderef{fred_m_na_stress_strain},
    \coderef{fred_m_na_stress_strain_residual}) — changing only the lines
    that need the new behaviour (here: 4 boundary-condition branch sites and
    the swelling term), and taking the application's own layer-state type
    instead of the shared \texttt{AxialLayerState}.
  \item Fork the residual-assembly loop the same way — see
    \texttt{FredMNaResiduals::computeResidualsMNa}
    (\coderef{fred_m_na_residuals_mna_nvi}) — reusing every protected
    \texttt{RodResiduals} helper (\texttt{layerPower}, \texttt{prepareLayer},
    \texttt{freezeThermalResiduals}, \ldots) and swapping in the new
    mechanics class in place of the inherited \texttt{m\_mech}.
  \item Register the forked loop's static \texttt{IDAResFn} instead of
    \texttt{RodResiduals::idaResidual} in the application's
    \texttt{getSolverResFn()} override.
  \item Keep any new app-specific state-machine logic (here:
    \texttt{FredMNaGapBehavior}, \coderef{fred_m_na_gap_behavior}) as its own
    class, called from the application's \texttt{afterAcceptedStep} override
    rather than the platform's IDA root-finder path.
\end{enumerate}
No file under \texttt{platform/} was touched; FRED-ROD/FRED-OX output is
therefore unaffected by construction, not merely by testing.

% ===========================================================================
\end{document}
% ===========================================================================
